\documentclass[script,pro]{debate}
\motion{对宠物的爱，爱的更是它 / 爱的更是我}
\stance{正方}{爱的更是它}
\meta{赛制}{中国科学技术大学新生辩论赛（四对四，正方先立论，反方先结辩）}
\meta{口径来源}{备赛文档 第十节 10.1 正方口径表。本文稿所有定义、判准、论点、数据均取自该表}
\begin{document}
\debatetitle
\begin{thesisbox}[全场一句话立论]钱和它的病相撞时多数人削的是自己；它病了老了不再回应时投入不降反升；明知结局仍然开始。冲突时让路的是我，所以爱的更是它。\end{thesisbox}
\begin{thesisbox}[全队三句（每个环节口径一致）]一、判准是冲突时优先谁：让谁受益、让谁让路。二、它的疼通过我的负担到达我的决定，渠道是我，起点和理由是它。三、被需要是我的需要；它需要我起床、我需要睡觉，起床的是我。\end{thesisbox}

\begin{speech}{1. 正一开篇立论稿}{正文 764 字 / 180 秒}
谢谢主席，问候在场各位。

我方的立场是：对宠物的爱，爱的更是它。

先说清三个词。爱，是对一个对象的深挚感情，以及这份感情带来的持续关切和行动。「更」，说明这是一道比较题，两边都在。我方承认养宠让人快乐、让人被需要，我方争的是分量。「它」，就是你家那一只，它有自己的需要，不随你的想法改变。

判准只有一条，也是最朴素的一条：冲突时优先谁。当它的需要和我的需要撞上，让路的是谁，这份爱就更是谁的。在这条判准下，我方要证明多数人让的是自己；对方要证明多数人让的是它。

我方第一点，回报归零的时候，投入最高。

二零一七年一项对照研究比较了二百三十八位猫狗主人，其中一百一十九位正在照护患慢性病或者终末期疾病的动物。照护组在负担、压力、抑郁和焦虑症状上全面更高，生活质量更差，每一项差异在统计上都极其显著。研究者的概括是，这种损耗在许多方面与照护一位患痴呆的家人相似。

你想想那是什么日子。去年新京报写过一个女孩，她的哈士奇从十三岁起失禁、瘫痪、痴呆，一晚上叫醒她六七次，她第二天照常上班，后来干脆辞了职。这个时候退路其实很便宜，少治、不治、送走，都在手边。可她没走，多数人也没走。

我方第二点，让路的是我。

再看一组更朴素的数字。二零二六年一项两千人的调查里，九成主人表示愿意为救宠物负债。而在过去一年，有三成七的人为了付兽医费，削减了自己的食品、水电和其他账单。

请评委记住这两个数字。削减自己的口粮去付它的药费。这就是「冲突时优先谁」最朴素的答案。

对方今天大概会说：你做这一切，图的是自己的满足感。可是数据摆在这里。当它病了、老了、不再回应你，它能提供的一切归零，主人的投入不但没有减少，反而在自己身上留下了可测量的伤害。回报最低的时候，投入最高。图回报，解释不了这件事。

所以我方全场都会请对方回答一个问题：在你方的标准下，一个饲主要做到什么，才算爱的更是它？说得出，我们逐条去比；说不出，这道题在你方口径下就不成立。

最后。没有人会为了爱自己，去主动选择一段注定在十几年后带来剧痛的关系。我们明知结局，\stress{仍然开始}。

所以我方坚定认为，对宠物的爱，是更爱它。谢谢。
\end{speech}
\begin{contingency}
\ifthen{反方把「爱」定义成“满足自我需要的情感投入”}{30 秒定义之争口径：这个定义让母爱、爱情全都“更是我”，命题不可证伪；请对方回答，按你方的标准有没有任何一种爱是更指向对方的。}
\ifthen{反方一开口就说“九成是意愿题”}{承认。然后指出三成七削口粮、三成七负债、四成四差点放弃但付了，全是行为题；行为题两边都不过半，人数相当，我方在钱上争的是方向。}
\ifthen{时间不够}{删新京报那个女孩的画面，只留“退路很便宜，可多数人没走”，其余不动。}
\end{contingency}

\begin{stage}{2. 正一被质询防守预案}{反四质询，120 秒双边计时}
\textbf{原则}：回答短，守住前提，不否认常识。全场绝对不能说的三句：“我什么都没得到”“被需要不是我的需要”“我的负担不决定它哪天走”。
\begin{qa}
\qarow{九成愿意负债，这是意愿题还是行为题？}{意愿题。行为题是三成七削了口粮、三成七负了债，也在同一份报告里。}
\qarow{同一份报告，四成六因成本延迟过就医，对吗？}{对。行为题两边都不过半，人数相当。我方在钱上争的是方向。}
\qarow{付不起就让它等，算不算让路？}{算。付不起的人一部分借钱削口粮，一部分让它等，谁也没证明多数。}
\qarow{被需要，是不是你的需要？}{是，我方承认。它需要我起床，我需要睡觉，起床的是我。}
\qarow{照护负担研究测的是你的负担，对吗？}{对。负担就是让出去的睡眠、时间、钱、健康。判准问谁让路，让路的是我。}
\qarow{同一作者说负担越高越考虑安乐死，撤不撤？}{不撤。它的疼通过我的负担到达我的决定，渠道是我，起点和理由是它。}
\qarow{安乐死的疼是主人感知的，对吗？}{对，它不会说话。观察的对象是它，兽医看的也是它。}
\qarow{你方是不是说付钱、停钱、撑着都是为它？}{不是。有得选时让它扛、拖着让它多疼、被告知后不改，这三样是为我。它们都存在，典型饲主做的是相反的事。}
\qarow{你方的比例数据是不是全在英美？}{多数是。中国数据证账单真实和机制存在，对方的中国比例数据分母也不对，双方对等。}
\qarow{你从养宠里得到了陪伴，对吗？}{对，我方第一句就承认了。「更」是比较，回报存在不等于回报更重。}
\end{qa}
\end{stage}

\begin{stage}{3. 正四质询反一问题链}{120 秒，双边计时}
\textbf{总目标}：拿到“存在更指向对方的爱”和“行为题两边都不过半”两个承认；逼反方说出“什么才算为它”。
\begin{chain}{链一：判准版定义杀}{逼反方给出可证伪的“为它”标准}
\q{母亲对孩子的爱，更是谁？}{答“孩子”：下一问；答“母亲”：那对方的标准下有没有任何一种爱更指向对方？}
\q{那按你方标准，饲主做什么才算爱的更是它？}{答不出：请评委记住，对方给不出标准；答得出（付得起就付、按它的指标放手、被告知后改）：下一问。}
\q{按它的指标放手，算不算为它？}{答“承认”：那近九成的狗在诊所以安乐死结束，理由八成七是它的疼；答“不承认”：那对方的标准是假的。}
\close{存在更指向对方的爱；按它的指标放手是为它；剩下的只是分量之争。}
\end{chain}
\begin{chain}{链二：行为题打平}{把“意愿 vs 行为”的框架拆掉}
\q{三成七削口粮、三成七负债，是不是行为题？}{答“是”：下一问；答“不是”：请念原文，“过去一年已经”。}
\q{行为题两边都不过半，对吗？}{答“对”：那“多数让它等”你方也没有；答“不对”：请给过半的数字。}
\q{急诊会拖的一成六，是不是比一般就医的四成六少？}{答“是”：冲突越重让它等的越少；答“那是假设题”：九成一也是假设题，同一标准。}
\close{钱的战场两边人数相当，对方要的“更”不在这里。}
\end{chain}
\begin{chain}{链三：渠道与起点}{防反方用负担中介接走我方的照护数据}
\q{负担中介那条路径，起点是不是它的疼痛？}{答“是”：下一问；答“否”：请念原文。}
\q{做出安乐死决定的理由，是不是它的疼？}{答“是”：那开关在它的指标上；答“那是事后填的”：你方的依恋主导也是同一张问卷填的。}
\cut{“好的，我理解为您没有否认。下一个问题。”}
\close{它的疼通过我的负担到达我的决定，渠道是我，起点和理由是它。}
\end{chain}
\end{stage}

\begin{speech}{4. 正二驳论稿}{正文 401 字 + 临场位 100 字 / 120 秒}
谢谢主席。

\reserve{此处留临场位，约 100 字：接反二驳论中最用力的一点。反二大概率打“负担高就考虑安乐死”或“被需要在病老期最高”。先复述对方原话，再用全队三句里对应的那一句接住，然后进入下面两个驳点。}
反方今天的立论，有一个自己都没绕过去的地方：他们要评委只看行为，可行为题里让自己的人，和让它等的人一样多。

反方第一个驳点，说九成愿意负债是嘴上说的，四成六延迟就医才是做的。好，只看做的。同一份报告里，三成七的人在过去一年已经削了自己的食品和水电，三成七已经真的负了债，四成四差点放弃但最后付了。这些也是行为。反方拿到的是“有人让它等过”，没有拿到“多数让它等”。而且请评委看方向：一般看病拖过的四成六，急诊说会拖的只有一成六。冲突越重，让它等的人越少。

反方第二个驳点更值得说。反方说病老之时它没有发言权，走或留跟着我的负担走。我方承认，我的负担会变。可反方引的那条路径，第一步写的是什么？是它的疼痛。它的疼通过我的负担到达我的决定，渠道是我，起点是它。真正做出的决定，八千九百多位犬主里八成七的理由是它的疼。你想想，同样一条路径，反方念的是中间那一步，我方念的是第一步和最后一步。反方把渠道当成了起点。

说到底，反方全场证明的是这份爱让我付了代价。代价落在谁身上，谁就让了路。判准问的正是这个。

\end{speech}
\begin{contingency}
\ifthen{反二打的是“被需要在病老期最高”}{临场位用全队第二句接：被需要是我的需要，它需要我起床、我需要睡觉，起床的是我。}
\ifthen{反二拿 61\% 被提示后不改}{那份问卷的分母是全体饲主，超重只有两到三成，七八成的宠物不胖；请对方给“被告知后不改”的比例。}
\ifthen{反二没打我方论点一}{把临场位的字数补到驳点二，加卢卡一晚起床六七次的画面。}
\end{contingency}

\begin{stage}{5. 正二对辩要点}{90 秒，单边计时}
\begin{points}{进攻点（每个 15–20 秒，说完就停）}
\item 行为题两边都不过半，反方要的“多数让它等”在哪里？
\item 负担中介的起点是它的疼。渠道是我，起点是它，反方把渠道当起点。
\item 依恋主导近七成，这个数字往哪边推？拖着，还是不忍它再疼？反方没有符号。
\item 在你方标准下，饲主做什么才算爱的更是它？这个问题反方还没答。
\item 反方说我的状况不随它的处境波动。二百三十八人的研究说，它病了，我全变了。
\end{points}
\begin{points}{备用点（回应对方进攻用）}
\item 得到是副产品，我方第一句就承认了。
\item 被需要是我的需要；起床的是我。
\item 付不起也是相撞，接受；那两边人数相当。
\item 八成四是假设题，我方不用它证多数；九成一也只说方向。
\item 观察的对象是它，兽医看的也是它。
\item 兽医被要求过无效治疗的比例不是饲主比例，反方自己承认了。
\item 六成一那份问卷的分母是全体饲主。
\item 对称的关系只能看谁让路，病老期让路的是我的睡眠和收入。
\end{points}
\begin{points}{对方最可能问的三个问题 → 回应}
\item “你承认九成是意愿题吗？” → 承认，行为题两边都不过半。
\item “负担越高越考虑安乐死怎么解释？” → 渠道是我，起点和理由是它。
\item “它的疼是你感知的吧？” → 对，它不会说话；观察的对象是它。
\end{points}
\end{stage}

\begin{stage}{6. 正三盘问问题链}{90 秒，单边计时}
\begin{chain}{链一：口径一致性检验}{让反二与反四在“按它的指标放手算不算为它”上出现分歧}
\q{问反二：按它的指标放手，算不算为它？}{答“算”：记住这句；答“不算”：那你方的标准是假的。}
\q{同一个问题问反四：您同意二辩的回答吗？}{答“同意”：口径锁死；答“不同意”：两人不一致。}
\q{那早放手就不能算它让路，对吗？}{答“对”：反方论点二只剩“拖”一条腿；答“不对”：两头都算，不可证伪。}
\q{“拖”的饲主比例，你方有吗？}{答“没有”：请评委记住；答“有”：请给。}
\close{反方承认按它的指标放手是为它，而“拖”没有饲主比例。}
\end{chain}
\begin{chain}{链二：对称关系}{打掉“我的状况不随它的处境波动”}
\q{问反四：它病了，主人的抑郁和负担变不变？}{答“变”：下一问；答“不变”：请解释二百三十八人的对照研究。}
\q{那我的状况是不是随它的处境波动？}{答“是”：不对称不存在；答“那是代价”：代价落在谁身上，谁就让了路。}
\q{问反二：病老期让出睡眠和收入的，是我还是它？}{答“我”：收；答“它”：请说它让出了什么。}
\close{关系是对称的，对称只能看谁让路，病老期让路的是我。}
\end{chain}
\textbf{盘问目标}：第一，反方在“早放手算不算为它”上要么口径不一，要么承认论点二只剩“拖”一条腿而没有饲主比例；第二，反方的“不对称”被自己引的数据否定。
\end{stage}

\begin{stage}{7. 正二、正四被盘问防守预案}{两人口径必须一致}
\begin{qaa}
\qaarow{九成愿负债是意愿题吧？}{是。行为题是三成七削口粮、三成七负债。}{是。行为题两边都不过半。}
\qaarow{付不起就让它等算不算让路？}{算。两边人数相当。}{算。我方在钱上争方向。}
\qaarow{被需要是不是你的需要？}{是。起床的是我。}{是。它需要我起床，我需要睡觉。}
\qaarow{负担越高越考虑安乐死，撤不撤？}{不撤。渠道是我，起点和理由是它。}{不撤。同一句。}
\qaarow{它的疼是你感知的吧？}{是。观察的对象是它。}{是。兽医看的也是它。}
\qaarow{饲主做什么才算为我？}{三样：有得选时让它扛、拖着让它多疼、被告知后不改。}{同上三样。}
\qaarow{你方的中国比例数据在哪？}{没有，我方承认；对方的分母也不对，双方对等。}{同上。}
\qaarow{临终补液是不是为了自己心安？}{有心安这一层，止痛的受益人是它。}{同上。}
\end{qaa}
两人赛前必须逐条对答一遍，答案可以换词，不能换意思。被逼到没准备过的问题时，统一说“我方立论的立场是”，不要现场发挥。
\end{stage}

\begin{speech}{8. 正三盘问小结稿}{正文 379 字 / 90 秒}
谢谢主席。

刚才的盘问里，我问了反方二辩和四辩同一个问题：按它的指标放手，算不算为它。请评委注意他们的回答。说算，那反方的论点二就只剩“拖着”这一条腿，而反方自己承认，拖着的饲主比例他们没有。说不算，那反方开场给的三条标准就是假的。

我还问了第二件事。反方全场都在说，它的处境跟着我的状况走，我的状况不随它的处境波动。可反方自己引的研究说的是什么？它病了，我的负担、抑郁、生活质量全变。二百三十八人的对照，一百五十七人的对照，两百七十七人的路径，三组数据说的都是我随它变。这个不对称不存在。

关系是对称的，那评委就只能看一件事：撞上的时候谁让路。老实说，钱的战场两边人数相当，我方不在那里争多数。可病老那一段，凌晨三点它叫，起床的是我；调动来了，辞职的是她。让出睡眠、收入、健康的是我，被照护的是它。

反方会说，代价不等于方向。请评委翻回判准：它从来没问过谁难受，它问的是撞上的时候谁让了路。到现在为止，双方的分歧只剩一个：代价落在我身上，算不算我让了路。我方认为算。

谢谢。
\end{speech}
\begin{contingency}
\ifthen{反二反四都答“算”}{直接推进：那反方论点二只剩“拖”，请给饲主比例。}
\ifthen{反三在我之后说“代价不等于方向”}{自由辩第一问就问：让出睡眠和收入的是谁。}
\end{contingency}

\begin{stage}{9. 自由辩战场设计}{240 秒}
\begin{battleground}{战场一：什么才算为它（前 90 秒，主战场）}
\begin{fields}
\field{目标}{让评委看到反方给不出可证伪的标准，或者给出的标准正好被我方数据满足。}
\field{开场问题}{按你方标准，饲主做什么才算爱的更是它？}
\field{对方答“付得起就付、按指标放手、被告知后改” → 追问}{近九成的狗在诊所以安乐死结束，理由八成七是它的疼，这算不算按指标放手？}
\field{对方答“给不出” → 追问}{那你方的命题不可证伪，请评委记住。}
\field{对方可能反攻 → 我方防守}{“你方的三条也过不了自己的标准。” → 拖着的饲主比例双方都没有；被告知后不改的比例反方拿不出；有得选时让它扛，两边人数相当。}
\field{收束语}{反方要的“更”，在他们自己的标准下也找不到。}
\end{fields}
\end{battleground}
\begin{battleground}{战场二：渠道与起点（中 90 秒）}
\begin{fields}
\field{目标}{把“负担中介”和“依恋主导”两个数字接回我方。}
\field{开场问题}{负担中介那条路径，起点是不是它的疼？}
\field{对方答“是” → 追问}{那渠道是我，起点是它，你方把渠道当成了起点。}
\field{对方答“同样的疼决定不同” → 追问}{那是决定有多难，做出的决定理由八成七是它的疼。}
\field{对方可能反攻 → 我方防守}{“依恋主导近七成。” → 这个数字往哪边推？拖着，还是不忍它再疼？没有符号。}
\field{收束语}{它的疼通过我的负担到达我的决定，渠道是我，起点和理由是它。}
\end{fields}
\end{battleground}
\begin{battleground}{战场三：对称关系与让路（最后 60 秒）}
\begin{fields}
\field{目标}{给评委一个画面收束全场。}
\field{开场问题}{它病了，主人的抑郁和负担变不变？}
\field{对方答“变” → 追问}{那关系是对称的，对称只能看谁让路；凌晨三点它叫，起床的是谁？}
\field{收束语}{代价落在谁身上，谁就让了路。让路的是我。}
\end{fields}
\end{battleground}
\begin{points}{必问（前两分钟内问完）}
\item 按你方标准，饲主做什么才算爱的更是它？（一辩稿抛出的那个问题）
\item 行为题两边都不过半，“多数让它等”在哪里？
\item 负担中介的起点是不是它的疼？
\item 依恋主导近七成，往哪边推？
\item 它病了，我的状况变不变？
\end{points}
\begin{points}{必答（15 秒以内正面回答）}
\item “九成是意愿题吧？” → 是。行为题是三成七、三成七、四成四，两边都不过半。
\item “负担越高越考虑安乐死。” → 渠道是我，起点和理由是它。
\item “被需要是你的需要吧？” → 是。起床的是我。
\item “它的疼是你感知的吧？” → 是，它不会说话。观察的对象是它。
\item “你方是不是说人怎么做都是为它？” → 不是。三样是为我，都存在，典型饲主做的是相反的事。
\end{points}
\textbf{时间分配与站位}：前 90 秒战场一，抛出必问一、二；中 90 秒战场二，处理对方进攻；后 60 秒战场三，不开新战场，留 20 秒备用。一辩守定义与判准，二辩接驳论过的点，三辩接盘问成果，四辩收束与价值。
\end{stage}

\begin{speech}{10. 正四结辩稿}{正文 736 字 / 210 秒，留 30–40 秒临场回应}
谢谢主席。

\reserve{此处留临场位，约 35 秒：真回应反四结辩。反四大概率会提前点破“老狗喂药”的画面，说“喂药是为了对自己有个交代”。先复述他这句原话，再说：交代完了，第二天她还是去喂了。然后进正文。不要背稿。}
全场打到现在，分歧其实只剩一句话：代价落在我身上，算不算我让了路。

反方今天最有力的一击，是那条路径：负担越高，越考虑安乐死。我方正面回答。这条路径的第一步写的是它的疼痛，第二步才是我的负担，第三步是我考虑。它的疼通过我的负担到达我的决定。渠道是我，起点是它。而真正做出的决定，八千九百多位犬主里，八成七的理由是它的疼。反方把渠道当成了起点，把“决定有多难”当成了“决定为了谁”。

反方的第二块证据，是同一份报告里四成六的人延迟过就医。我方承认，人数相当。可同一份报告里三成七的人削了自己的食品水电，三成七负了债。反方拿到的是“有人让它等过”，从来没有拿到“多数让它等”。说白了，反方全场证明的是这份爱让我付了代价。可代价落在谁身上，谁就让了路。

我方的三个论点还立在这里。钱和它的病相撞时，削自己口粮的人在那里。它病了老了不再回应时，照护它的人在自己身上留下了可测量的伤害，还在照护。到终点，多数人按它的痛，不按自己的不舍做决定。请评委看一眼这三个时刻的共同点：让出去的，每一次都是我的东西。

我方一辩抛出的那个问题，反方到现在也没有正面回答：按他们的标准，饲主要做到什么，才算爱的更是它。

我想说一件小事。

去年新京报写过一个女孩。她的哈士奇十三岁那年开始失禁，然后瘫痪，然后痴呆。它不再在门口等她，不再认得她。它一晚上叫醒她六七次，她起来，收拾，躺下，再起来，第二天照常上班。后来她辞了职，每三个小时回一趟家。

在那段日子里，这段关系已经不给她任何东西了。没有陪伴，没有回应，连一个认得她的眼神都没有。

她还在做。她不是特例，二百三十八人的对照研究里，这样的人是一整组。

到它连水都咽不下去的那天，她签了字。签字的时候，它少疼了，她的日子空了。

各位，今天我们争的不是谁在这份感情里得到得更多。我们争的是，一个人有没有能力去爱一个永远不会说「我也爱你」的生命。

我方认为有。而且我方认为，恰恰是在它什么都给不了你的那一段日子里，这份爱才露出了它本来的方向。

它一直朝着它。

谢谢。
\end{speech}
\begin{contingency}
\ifthen{反四已经把“老狗喂药”的画面提前点破（封路）}{不换画面，换讲法：承认反四说对了一半，“她确实在对自己交代”，然后加一句“可交代完了，第二天她还是去喂了”。}
\ifthen{反四把判准整个改成“决定变量在我”}{临场位先用 10 秒说：判准两问，让谁受益、让谁让路，决定由谁做不是判准。再进正文。}
\ifthen{我方某条数据被打穿}{不硬撑，收缩到二百三十八人的对照与卢卡的画面，用判准解释为什么仍然占优。}
\end{contingency}
\end{document}
